\documentclass[11pt]{article}
\usepackage[paperwidth=11in,paperheight=8.5in,margin=0.25in]{geometry}
\usepackage[T1]{fontenc}
\usepackage[utf8]{inputenc}
\usepackage[english]{babel}
\usepackage{lmodern}
\usepackage{microtype}
\usepackage{amsmath,amssymb}
\usepackage{xcolor}
\usepackage{tikz}
\usetikzlibrary{calc}
\hyphenpenalty=8000
\exhyphenpenalty=8000
\emergencystretch=2em
\definecolor{QCBlue}{HTML}{1F4E79}
\definecolor{QCGreen}{HTML}{2E6B5F}
\definecolor{QCOrange}{HTML}{B65B2A}
\definecolor{QCPurple}{HTML}{4B3F72}
\definecolor{QCLight}{HTML}{F4F6F8}
\definecolor{QCBorder}{HTML}{D0D6DC}
\definecolor{QCText}{HTML}{111111}
\newlength{\Margin} \setlength{\Margin}{0.25in}
\newlength{\Gap} \setlength{\Gap}{0.1in}
\newlength{\TitleH} \setlength{\TitleH}{0.8in}
\newlength{\HeaderH} \setlength{\HeaderH}{0.38in}
\newlength{\Pad} \setlength{\Pad}{0.1in}
\newlength{\UsableW}
\newlength{\CellW}
\newlength{\CellH}
\setlength{\UsableW}{\dimexpr\paperwidth-2\Margin\relax}
\setlength{\CellW}{5.2in}
\setlength{\CellH}{3.5in}
\newcommand{\QuadBox}[4]{%
  \node[anchor=north west, draw=QCBorder, line width=0.6pt, rounded corners=10pt, fill=white, minimum width=\CellW, minimum height=\CellH, inner sep=0pt] (B) at #1 {};
  \fill[#2, rounded corners=10pt] ([xshift=0pt,yshift=0pt]B.north west) rectangle ([xshift=0pt,yshift=-\HeaderH]B.north east);
  \node[anchor=west, text=white, font=\bfseries] at ([xshift=0.18in,yshift=-0.22in]B.north west) {#3};
  \node[anchor=north west, text=QCText] at ([xshift=\Pad,yshift=-\HeaderH-\Pad]B.north west) {\begin{minipage}[t]{\dimexpr\CellW-2\Pad\relax}\raggedright\small #4\end{minipage}};
}
\pagestyle{empty}
\begin{document}
\begin{tikzpicture}[remember picture,overlay]
\coordinate (NW) at ([xshift=\Margin,yshift=-\Margin]current page.north west);
\node[anchor=north west, draw=QCBorder, line width=0.6pt, rounded corners=12pt, fill=QCLight, minimum width=\UsableW, minimum height=\TitleH, inner sep=0pt] (T) at (NW) {};
\node[anchor=north west] at ([xshift=0.22in,yshift=-0.22in]T.north west) {\begin{minipage}[t]{\dimexpr\UsableW-0.44in\relax}{\Huge\bfseries \textcolor{QCBlue}{Electromagnetic Radiation Formula Sheet}}\par\vspace{2pt}{\normalsize \textbf{Diego Juarez}\hfill \textbf{Computational Electrodynamics}\hfill \textbf{\today}}\end{minipage}};
\coordinate (GridNW) at ([yshift=-0.15in]T.south west);
\coordinate (TL) at (GridNW);
\coordinate (TR) at ([xshift=\CellW+\Gap]GridNW);
\coordinate (BL) at ([yshift=-\CellH-\Gap]GridNW);
\coordinate (BR) at ([xshift=\CellW+\Gap,yshift=-\CellH-\Gap]GridNW);
\QuadBox{(TL)}{QCBlue}{1) Retarded Potentials}{
Lorenz gauge:
\[
\Phi(\mathbf r,t)=\frac{1}{4\pi\varepsilon_0}\int \frac{\rho(\mathbf r',t_r)}{|\mathbf r-\mathbf r'|}\,d^3r'
\]
\[
\mathbf A(\mathbf r,t)=\frac{\mu_0}{4\pi}\int \frac{\mathbf J(\mathbf r',t_r)}{|\mathbf r-\mathbf r'|}\,d^3r'
\]
Retarded time:
\[
t_r=t-\frac{|\mathbf r-\mathbf r'|}{c}
\]
}
\QuadBox{(TR)}{QCGreen}{2) Radiation Zone}{
Near field:
\[
\sim \frac{1}{r^3}
\]
Induction field:
\[
\sim \frac{1}{r^2}
\]
Radiation field:
\[
\sim \frac{1}{r}
\]
Far zone:
\[
\mathbf B=\frac{1}{c}\hat{\mathbf r}\times \mathbf E
\]
\[
\mathbf E\perp \mathbf B \perp \hat{\mathbf r}
\]
Poynting vector:
\[
\mathbf S=\frac{1}{\mu_0}\mathbf E\times \mathbf B
\]
}
\QuadBox{(BL)}{QCOrange}{3) Oscillating Dipole}{
Dipole moment:
\[
\mathbf p(t)=p_0\cos(\omega t)\,\hat{\mathbf z}
\]
Far-zone field amplitudes:
\[
E_\theta \propto \frac{\omega^2 p_0 \sin\theta}{4\pi\varepsilon_0 c^2 r}
\]
\[
B_\phi=\frac{E_\theta}{c}
\]
Angular pattern:
\[
I(\theta)\propto \sin^2\theta
\]
Maximum at $\theta=\pi/2$, zero on the dipole axis.
}
\QuadBox{(BR)}{QCPurple}{4) Power and Checks}{
Average radiated dipole power:
\[
\langle P\rangle \propto \frac{\omega^4 p_0^2}{\varepsilon_0 c^3}
\]
Nonrelativistic accelerating charge:

Radiation requires acceleration, not uniform motion.

\textbf{Checks}

1. Radiation fields are transverse.

2. Only $1/r$ terms carry energy to infinity.

3. Symmetry can kill dipole radiation and elevate quadrupole terms.
}
\end{tikzpicture}
\end{document}
